% Preâmbulo compartilhado dos PDFs do livro — pensado para qualidade de impressão.
\documentclass[12pt, a4paper, oneside, open=any]{scrbook}

\usepackage{fontspec}
\usepackage{polyglossia}
\setmainlanguage{brazil}
% DejaVu Serif tem cobertura Unicode ampla (símbolos como ⊂ usados no capítulo
% de arquitetura), diferente da Latin Modern padrão.
\setmainfont{DejaVu Serif}
\setsansfont{DejaVu Sans}
\setmonofont{DejaVu Sans Mono}

\usepackage{amsmath}
\usepackage{amssymb}
\usepackage{microtype}
\usepackage{geometry}
\geometry{left=2.8cm, right=2.2cm, top=2.8cm, bottom=2.8cm}

\usepackage{booktabs}
\usepackage{longtable}
\usepackage{array}
\usepackage{calc}
\usepackage{enumitem}
\usepackage{xcolor}
\usepackage{tcolorbox}
\tcbuselibrary{skins, breakable}

\usepackage{graphicx}
\setkeys{Gin}{width=0.85\linewidth, keepaspectratio}

\usepackage[hidelinks]{hyperref}
\usepackage{fancyhdr}
% (scrbook/KOMA-Script já provê estilização de seção nativa; titlesec conflita
% com a classe e foi removido de propósito)

% Cabeçalho corrente: nome do livro à esquerda, capítulo à direita
\pagestyle{fancy}
\fancyhf{}
\fancyhead[LE]{\small\nouppercase{\leftmark}}
\fancyhead[RO]{\small\nouppercase{\rightmark}}
\fancyfoot[C]{\thepage}
\renewcommand{\headrulewidth}{0.4pt}

% Numeração e estilo de seção herdados do KOMA-Script (scrbook já cuida bem disso)
\setcounter{secnumdepth}{3}
\setcounter{tocdepth}{2}

% Caixa de destaque para "Figura pendente" (usada pelo filtro Lua)
\newtcolorbox{figurapendente}{
  colback=yellow!8, colframe=yellow!60!black,
  title=Figura pendente, fonttitle=\bfseries,
  breakable
}

% Caixa de destaque para analogias/exemplos didáticos (opcional, uso futuro)
\newtcolorbox{analogiabox}{
  colback=blue!4, colframe=blue!40!black,
  breakable
}
